% !TeX root = ../深度学习与大语言模型技术原理.tex
% 由主文档通过 \input 载入；请编译主文档。
\setcounter{chapter}{0} % 保留原书章号。
\chapter{绪论：从可学习的函数到可验证的系统}
\label{ch:introduction}

\section{研究对象与问题层次}

\emph{深度学习（deep learning）}研究如何用多层可微变换，从数据中学习完成任务的表示与预测规则。给定参数为 $\theta$ 的函数 $f_\theta$，学习过程并不是给每个输入编写一条规则，而是设计模型族、数据分布和目标函数，再用优化算法选择参数。\emph{大语言模型（large language model, LLM）}是这一思想在离散序列上的重要实例：它根据已有词元预测下一个词元，并通过不断重复这个过程生成文本、代码或结构化动作。

预测下一个词元是训练接口，却不足以说明整个系统。模型是否能检索正确证据、遵守工具参数约束、在有限显存中服务并发请求，以及是否能被可靠评估，分别涉及信息检索、程序接口、计算系统与统计实验。因而，理解一项技术至少要回答四个相互关联的问题：它修改了什么对象，优化了什么目标，付出了什么成本，以及在什么条件下成立。例如，低秩适配改变可训练参数的表示；检索增强生成改变模型可见的信息；分页键值缓存改变存储管理；约束解码改变允许生成的词元集合。这些方法作用于不同层次，可以组合，却不能相互替代。

本书以同级目录中的《技术清单.txt》为主题来源。原清单的第1至71节提供主要技术范围，第0节补充世界模型和稀疏混合架构，第72至82节中的项目与学习建议转化为综合实验和阅读路径。重复出现的技术合并讲解，清单对应关系见附录~\ref{app:map}。本书不把清单中的年度热度、求职收益和岗位竞争力排序当成已验证的事实，也不把软件品牌当作独立数学理论。

\section{阅读路线与知识依赖}

全书先建立线性代数、概率和优化语言，再进入分词、语言建模、Transformer、专家混合、数据与分布式训练。后训练部分依次讨论监督微调、偏好建模、策略优化、可验证奖励和蒸馏，使每个目标函数都能追溯到清楚的数据来源。随后介绍检索增强生成与智能体，将模型放入有外部信息和环境反馈的系统中。评测与安全用于检验这些系统的行为；高效推理、模型压缩、长上下文、多模态和底层工程则解释如何实现与扩展它们。

初次阅读可以沿章序建立主线。已有训练基础的读者可以从第~\ref{ch:sft} 章开始学习后训练，再回查第~\ref{ch:probability} 章的概率记号。偏重应用系统时，先阅读第~\ref{ch:retrieval} 至第~\ref{ch:agentlearning} 章，并结合第~\ref{ch:evaluation} 章设计评测。偏重训练与部署时，应把第~\ref{ch:stability}、\ref{ch:distributed} 章与第~\ref{ch:inference} 至第~\ref{ch:speculation} 章连起来阅读。

这里给出的推导是方法理解的基础，实验效果则必须在指定模型、数据和预算下重新检验。论文中的一种成功实现，不是所有模型都满足的最优性定理。软件功能也依赖版本、硬件、内核和模型结构；本书只介绍工具职责和稳定概念，实际接口应以所固定版本的官方文档为准。

\section{统一记号与成本口径}

全书默认单个表示向量为列向量，线性层写为 $y=Wx$，其中 $W\in\mathbb R^{d_{\mathrm{out}}\times d_{\mathrm{in}}}$。在批量注意力中，把词元表示按行堆叠为 $X\in\mathbb R^{n\times d}$，于是同一个列向量线性变换写为 $XW^\top$。采用不同堆叠方式时会明确说明，不能忽略转置混用这两种约定。

用 $n$ 表示序列长度，$d$ 表示隐藏维度，$L$ 表示层数，$B$ 表示批量大小，$V$ 表示词表大小，$h_q$ 与 $h_{kv}$ 分别表示查询头数和键值头数，$d_h$ 表示每头维度。模型参数总数记为 $P$，可训练参数数目记为 $P_{\mathrm{train}}$。$\mathcal D$ 表示数据集或明确指出的数据分布，$\theta$ 表示模型参数；奖励模型参数另记为 $\phi$，价值模型参数另记为 $\psi$。

随机变量用大写字母，取值通常用小写字母；上下文中的矩阵 $X$ 与随机变量 $X$ 依对象说明区分。$p_\theta$ 表示条件概率模型；讨论强化学习时，同一个生成模型写为策略 $\pi_\theta$。$\log$ 默认是自然对数，信息量单位为奈特。复杂度中的 $O(\cdot)$ 隐去常数与未讨论的维度，但不能把运算量直接等同于实测时间。一次乘加按两个浮点运算计数时，会明确采用这个口径。存储的 GiB 为 $2^{30}$ 字节，GB 为 $10^9$ 字节。

